% Chapter 53: From One Electron to the Periodic Table
\chapter{From One Electron to the Periodic Table}

Chapter~52 toured hydrogen's showroom: one Coulomb-bound electron occupies discrete states labeled by four quantum numbers. Look around: your iron spoon has twenty-six electrons, a gold ring seventy-nine, atmospheric neon ten. If multiple electrons merely enlarged hydrogen, more-electron elements should resemble one another increasingly, making a monotonous world. Instead sodium explodes in water while neighboring neon refuses every reaction; carbon builds life, silicon chips. A hundred diverse elements grow from identical quantum rules. Two new rules, indistinguishability and one electron per quantum state, will organize them into the periodic table.

\step{A Counterintuitive Opening}

Press your palm flat against a table. It does not budge, unremarkably until Chapter~51's scattering returns: nuclei occupy only a trillionth of atomic volume, leaving a few electron clouds in vast emptiness. Why do hand and table, both almost void, encounter an invisible solid wall instead of passing through?

Do not rush to charge repulsion. Both are neutral, so long-range Coulomb forces should vanish, or polarization should attract slightly, never make an immovable wall. Yet a table supports a person. This wall also sets water's incompressibility, metal hardness, and dead stars' resistance to gravity. No textbook names this world-supporting new force because it is not a new force at all.

Compare neon tubes, available at hardware stores and glowing when powered, with neighboring sodium. Neon never reacts with anything; sodium, element eleven beside neon's ten, hisses, burns, or explodes when a piece enters water. Why does one extra electron turn utter laziness into extreme temper? Similar personalities recur every few elements, two, eight, eight, eighteen, a code old physics never explained. It will speak by this chapter's end.

\step{Make a Guess First}

Write three predictions before proceeding.

First: what blocks your palm, (a) colliding nuclei, (b) electrostatic electron repulsion, or (c) a charge-independent new mechanism? If (b), why do positive nuclei and negative clouds in neutral atoms not cancel attractions and repulsions, leaving repulsion specifically?

Second: as nuclear charge and electron count increase, does atomic size grow, shrink, or remain similar while varying periodically? Does the energy to remove the outermost electron rise monotonically or fluctuate?

Third and crucially: why period lengths \(2,8,8,18,18,\dots\)? Rewrite them as \(2\times1\), \(2\times4\), \(2\times4\), \(2\times9\), \(2\times9\). See a pattern? Where does the factor \(2\) originate?

Bookmark the answers; each section will revisit one.

\step{Hands-on Experiment}

\begin{experiment}[Two Walls Inside a Syringe]
\textbf{Materials}: a clean plastic syringe without a needle, water, optional table salt, and an optional gas stove or candle with adult supervision.

% BEGIN TRANSLATOR SAFETY NOTE
\textbf{Translator safety note:} Do not force a blocked syringe or apply body weight as the original steps suggest. Use a recorded demonstration, or an instructor-approved needle-free activity with only gentle hand pressure, stopping at resistance and keeping the outlet away from faces. Overpressure can cause equipment failure; see \href{https://www.hse.gov.uk/pressure-systems/about.htm}{HSE pressure-equipment safety guidance}.
% END TRANSLATOR SAFETY NOTE

\textbf{Step one}: remove the needle, half-fill with air, tightly block the outlet with a finger, and push hard. Air visibly compresses and rebounds upon release. Record the approximate volume reduction.

\textbf{Step two}: fill completely with bubble-free water, block again, and push harder. The plunger barely moves, not because of friction, since it moves smoothly when gently jiggled, but because water resists compression. Even applying your whole body's weight makes no visible volume change.

% BEGIN TRANSLATOR SAFETY NOTE
\textbf{Translator safety note:} Do not perform the optional stove/candle step at home or hold copper wire in a flame. Use a recorded demonstration or an instructor-run flame test with suitable equipment and splash goggles; never add flammable solvents. See \href{https://institute.acs.org/acs-center/lab-safety/education-training/safer-experiments/flame-test.html}{ACS flame-test guidance}.
% END TRANSLATOR SAFETY NOTE

\textbf{Step three, optional}: sprinkle salt into the stove or candle flame for bright yellow; insert cleaned copper wire for green. Each element announces its name.

\textbf{Record and think}: both air and water contain almost-empty atoms or molecules. Why can one halve its volume while the other resists even a thousandth? Where inside the syringe is the wall resisting your finger? Explain particularly why electrically neutral water molecules have net repulsion if electrostatics is responsible.
\end{experiment}

We will quantify this: one-percent water compression needs about \(2\times10^{7}\)~Pa, a two-hundred-kilogram person on each square centimeter. Your strongest syringe push falls two orders short. Salt's yellow light signs sodium's temperamental outer electron; later you will understand why yellow.

\step{The Old Model Runs into Trouble}

Our best old tools are Coulomb's law and Chapter~52's quantum hydrogen. Apply both seriously to table and syringe and watch them fail.

First, neutral atoms should have no wall. At distance their charge attractions and repulsions cancel, leaving weak polarization attraction, the very attraction condensing water. Overlapping clouds should initially attract more. Even further compression has no classical prohibition against opposite-charge clouds interpenetrating. Matter should pass like smoke or flatten like clay, yet syringe water is rock-hard. This is qualitative failure, not numerical error.

Second, multiple electrons should not produce periodicity. With nuclear charge \(+Ze\), all \(Z\) electrons seek the lowest state, Chapter~52's smallest cloud of radius roughly \(a_0/Z\), where \(a_0\) is hydrogen's roughly \(5.3\times10^{-11}\)~m radius. Uranium, element ninety-two, should then be ninety-two times smaller than hydrogen, below a picometer. Its measured radius is about \(1.5\times10^{-10}\)~m, roughly three times hydrogen, over a hundredfold discrepancy. Likewise, inward-packed electrons should become monotonically harder to remove as \(Z\) rises. But Mendeleev found periodic personalities in 1869: lithium, sodium, potassium form one family; fluorine, chlorine, bromine another; noble gases end each with collective idleness. A monotonic model cannot generate a periodic table.

Third, is two electrons merely twice one electron? Classical billiard balls retain labels and trajectories. Every measured electron mass, charge, and spin is identical within resolution; no unusual electron has been found. Without identity cards, are electron one left/electron two right and the reversed assignment one state or two? This apparent word game is the periodic table's birth certificate.

\step{Inventing New Concepts}

Take the word game seriously, as physicists did, and invent three things.

First, \textbf{identical particles}. Promote observed indistinguishability to an accounting rule: exchanging electrons creates no new state. Stronger than unmeasurable differences, it denies differences in principle, immediately changing counting. Black and white socks permit four ways of drawing; labeled balls exchanged give two arrangements; exchanged electrons represent the same reality, so double-counting is wrong. Like bank money, swapping two hundred-unit deposits changes no balance. Unlike identical-looking billiard balls, electrons cannot be continuously tracked and tagged without disturbing their states.

Second, \textbf{two exchange characters}. Identical particles have two quantum accounting possibilities: exchange leaves the state unchanged, defining \textbf{bosons}, or changes its overall sign, defining \textbf{fermions}. An overall sign seems harmless because probabilities ignore it, yet consequences transform everything. Nature strictly pairs integer spin, photons and helium-4 atoms, with bosons, and half-integer spin, electrons, protons, neutrons, with fermions. Accept spin--statistics here as experimental fact; its deeper relativistic field-theory reason belongs to Chapter~49. Electrons are spin-one-half fermions, the periodic table's seed.

Third, our protagonist, \textbf{Pauli exclusion}. Suppose two electrons share exactly one quantum state. Exchanging them changes nothing, yet fermions require a sign reversal. Only zero equals its own negative, so that hypothetical state does not exist. Atomically, \textbf{no two electrons in one atom can share all four quantum numbers}. Not difficult or expensive: the account has no such entry. Notice that exclusion does not forbid shared position, since clouds overlap; it forbids a shared state, one electron per room addressed by Chapter~52's \(n\), \(l\), \(m_l\), \(m_s\).

Now multiple-electron structure has a skeleton. Electrons fill lowest-energy rooms first. Principal \(n\) sets floors, angular \(l\) room types, \(l=0,1,2,3\) denoting s, p, d, f subshells; magnetic \(m_l\) sets orientation and spin \(m_s\) has two choices. Floor \(n\) therefore has \(2n^2\) rooms: two, then eight. The code \(2\times1\), \(2\times4\), \(2\times9\) gets its \(2\) from spin's orientations. Period lengths \(2\) and \(8\) fall directly from counting quantum numbers.

Real energies need correction: electrons screen one another. Inner clouds conceal some nuclear charge, so outer electrons feel not \(+Ze\) but reduced \textbf{effective nuclear charge}. Screening depends on shape. s clouds penetrate near the nucleus, encounter less screening, and have lower energies; d and f largely remain outside inner electrons. Energies therefore cut across floors: fourth-floor 4s falls below third-floor 3d, producing 1s, 2s, 2p, 3s, 3p, 4s, 3d, 4p. In \(2,8,8,18\), the second \(8\) and first \(18\) result from this queue-jumping.

Within a subshell, \textbf{Hund's rule} fills equivalent orbitals singly with parallel spins before pairing. Two physical reasons: separate orbitals reduce Coulomb repulsion through spatial separation; parallel spins further lower energy through exchange, quantum accounting rather than a new force. Hund is not a fundamental law but electrostatics plus identical-particle rules, selecting the lowest-energy option among nearly equivalent fillings.

Now helium-3 and helium-4 supply a twin counterexample. With two electrons each, their chemistry is nearly identical, suggesting similar low-temperature behavior. Wrong: helium-4 has two protons, two neutrons, and two electrons, six fermions whose even total makes a composite boson. Below \(2.17\)~K it becomes a viscosity-free superfluid that climbs container walls. Helium-3 lacks one neutron, retaining fermionic character with five, and needs roughly \(2.5\)~mK, nearly a thousand times colder, for superfluidity through different pairing. One neutron barely matters chemically but makes two physical worlds, exclusive-seat fermions versus many-in-one-state bosons. Intuition ignores constituent parity; nature counts.

\step{The Model in One Sentence}

\begin{oneline}
Identical fermionic electrons occupy quantum states one apiece, filling lowest-energy rooms before moving upstairs. Screening and penetration adjust the filling rhythm into the periodic table, while outermost electron counts determine chemical personality.
\end{oneline}

Like a theater, sell one seat each, best first, opening rows successively. Three differences: seats are four-number states, not spatial cells, and different-state electrons can share position; energy and Hund's rule determine seating, not preference or arrival order, with alternatives forbidden or more costly; and full noble-gas shells resist addition because opening another floor exceeds chemistry's budget, not because occupants are lazy.

\step{The Mathematical Translator}

Count. Given \(n\), \(l\) takes \(0,1,\dots,n-1\); given \(l\), \(m_l\) takes \(-l,\dots,0,\dots,+l\), totaling \(2l+1\). Double each orbital for spin. Floor \(n\) holds
\[
N_n=\sum_{l=0}^{n-1}2(2l+1)=2n^2 .
\]
Four questions: it counts maximum shell occupancy; \(2l+1\) counts orientations, \(2\) spins, and the sum room types, elementary arithmetic-series mathematics. It applies when four quantum numbers approximately describe each electron, the independent-electron approximation scrutinized below. The counting never errs, but strong correlations undermine the independent-orbital picture. Check \(n=1\), \(N_1=2\), hydrogen and helium; \(n=2\), \(N_2=8\), the second period; \(n=3\) gives \(18\), yet period three has \(8\). No miscount: 4s delays 3d's ten rooms until floor four opens. A good formula's apparent mismatch locates hidden physics, here screening.

Estimate effective nuclear charge by treating an outer electron in the nucleus-plus-inner-cloud Coulomb field. Borrow Chapter~52's hydrogen ionization estimate:
\[
I\approx 13.6\ \text{eV}\times\frac{Z_{\mathrm{eff}}^2}{n^2},
\qquad Z_{\mathrm{eff}}=Z-\sigma ,
\]
where \(\sigma\) counts screened charges. Check \(\sigma=0\), no inner electrons, hydrogen, recovering \(13.6\)~eV. For \(\sigma=Z-1\), complete screening leaves \(Z_{\mathrm{eff}}=1\), like excited hydrogen; real atoms lie between. Sodium has \(Z=11\), outer 3s with \(n=3\), and measured \(I=5.14\)~eV. Invert:
\[
Z_{\mathrm{eff}}=n\sqrt{\frac{I}{13.6\ \text{eV}}}=3\times\sqrt{\frac{5.14}{13.6}}\approx 1.8,
\qquad \sigma\approx 11-1.8\approx 9.2 .
\]
Ten inner electrons screen about \(9.2\) nuclear charges, almost but not completely, evidence of 3s penetration. Lithium, \(Z=3\), outer 2s and \(I=5.39\)~eV, gives \(Z_{\mathrm{eff}}\approx1.3\), \(\sigma\approx1.7\): two 1s electrons screen \(1.7\) charges. Two data lines pin down inner shade and outer glimpses of nuclear sunlight.

Screening creates two worlds in one atom. Uranium's \(Z=92\) innermost 1s electrons nearly feel all \(92\) charges, binding at \(13.6\times92^2\approx1.15\times10^5\)~eV, matching measured K-shell \(1.16\times10^5\)~eV. Outermost electrons feel only a few effective charges and ionize near \(6\)~eV, a nearly twenty-thousandfold difference. Chemistry involves only these outer eV-scale electrons, explaining its energy scale and the periodic table's reach.

\begin{figure}[htbp]
\centering
\begin{tikzpicture}[>=Stealth,scale=0.62]
  % Energy levels
  \foreach \y/\lab in {0/{1s}, 1/{2s}, 2/{2p}, 3/{3s}, 4/{3p}, 5/{4s}, 6.6/{3d}, 7.6/{4p}}{
    \draw[thick] (0,\y) -- (3,\y) node[right=2pt]{\lab};
  }
  \node[left] at (0,3.8) {Energy};
  \draw[->] (-0.7,-0.3) -- (-0.7,7.9);
  % Filling-order arrows
  \draw[->,dashed,blue!60!black] (1.5,0) -- (1.5,0.9);
  \draw[->,dashed,blue!60!black] (1.5,1.1) -- (1.5,1.9);
  \draw[->,dashed,blue!60!black] (1.5,2.1) -- (1.5,2.9);
  \draw[->,dashed,blue!60!black] (1.5,3.1) -- (1.5,3.9);
  \draw[->,dashed,blue!60!black] (1.5,4.1) -- (1.5,4.9);
  \draw[->,dashed,blue!60!black] (1.5,5.1) -- (1.5,6.5);
  \draw[->,dashed,blue!60!black] (1.5,6.7) -- (1.5,7.5);
  % Electron occupancy
  \foreach \y/\num in {0/2, 1/2, 2/6, 3/2}{
    \node at (1.5,\y+0.25) {\scriptsize \num electrons};
  }
  \node[align=left] at (5.9,2.6) {\scriptsize Through 3p:\\ \scriptsize eighteen electrons,\\ \scriptsize argon's full shell};
\end{tikzpicture}
\caption{Multiple-electron ticket sales: energy levels, not to scale, and dashed filling arrows. Penetration places 4s below 3d, accounting for different third- and fourth-period lengths.}
\end{figure}

\begin{mathbridge}
\textbf{Exchange symmetry in formulas.} Write a two-particle state \(\psi(\mathbf x_1,\mathbf x_2)\). Identical-particle physics is unchanged by relabeling: \(\psi(\mathbf x_2,\mathbf x_1)=e^{i\alpha}\psi(\mathbf x_1,\mathbf x_2)\). Two exchanges restore the original, so \(e^{2i\alpha}=1\), permitting \(e^{i\alpha}=\pm1\): \(+1\) bosons, \(-1\) fermions. If two fermions occupy \(\phi\), their \(\psi=\phi(\mathbf x_1)\phi(\mathbf x_2)\) does not change sign, contradicting fermionic exchange, hence \(\psi=0\). An \(N\)-electron state must be fully antisymmetric. Its Slater-determinant form visibly vanishes with identical rows: one-seat exclusion is the algebra of equal determinant rows.

\textbf{Degeneracy pressure: exclusion made mechanical.} Put \(N\) electrons into volume \(V\). Exclusion fills progressively higher momentum states up to the Fermi energy
\[
E_F=\frac{\hbar^2}{2m_e}\left(3\pi^2\frac{N}{V}\right)^{2/3}.
\]
As density vanishes, \(E_F\) vanishes, no crowding, no repulsion; increasing density raises \(E_F\) as power \(2/3\). Crucially, even at absolute zero, electrons retain enormous kinetic energy, and compressing \(V\) costs more. The energy-volume derivative gives pressure. This \textbf{degeneracy pressure} supports tables and white dwarfs, not a new interparticle force but the mechanical bill for prohibited state sharing.
\end{mathbridge}

\begin{challenge}
\textbf{Why one electron per state?} Spin--statistics needs relativistic quantum field theory. In three dimensions, microscopic causality, spacelike measurements not affecting one another, plus positive energy require half-integer fields to use fermionic anticommutation and integer fields bosonic commutation. Exclusion is not a detachable empirical appendix; removing it collapses consistency. Century-long searches, some bounding violations below \(10^{-30}\), remain empty without surprising physicists.

\textbf{Exchange energy: Hund's quantitative origin.} Electrons in orbitals \(a\) and \(b\) acquire an antisymmetric wavefunction forbidding coincidence, the Fermi hole. Parallel spins thereby reduce Coulomb repulsion. An orbital-overlap exchange integral appears in energy corrections, lowering parallel-spin states. Read the cause accurately: not magnetic attraction, but identical-particle accounting changes electrostatic energy. This classically unmatched exchange underlies both Hund and Chapter~54's ferromagnetism.
\end{challenge}

\step{The Model's Limits}

Separate the customary three piles.

\textbf{Experimental facts}: electron mass, charge, and spin agree within measurement precision; spectra, ionization, and chemistry recur with periods \(2,8,8,18,18,32\); closed noble shells are inert, single outer alkali electrons active; helium isotopes differ radically in superfluidity; no confirmed Pauli violation exists. These are interpretation-independent.

\textbf{Mathematical model}: independent electrons occupy four-number orbitals in a nuclear-plus-average-electron field, subject to exclusion. This explains periodic structure, ionization trends, valence, and approximate energies remarkably well. It fails when correlations defeat mean fields, as in some transition-metal oxides and heavy-fermion materials, requiring genuine many-body theory. Precise energies, such as a third ionization-energy significant digit, likewise need many-body methods. Screening and penetration are themselves approximate model language, not fundamental mechanisms.

\textbf{Physical explanation}: relativistic spin--statistics answers why two statistics and why spin determines them; this is not an unresolved interpretation dispute. But premises include three spatial dimensions and Lorentz invariance. Two-dimensional exchanges can have arbitrary anyon phases, experimentally observed. Even the fermion/boson division has boundaries.

Three common misuses: Pauli is not a new inverse-square force but a state constraint whose energy-volume consequences yield pressure, analogous to centripetal not being a new force. Orbitals are still not planetary paths: spherical s clouds do not circulate, and penetration describes probability shape, not occasional shortcuts. Finally, crowd-shy fermions and sociable bosons have no psychology. Photons enter occupied modes not through attraction but because bosonic probability weights favor modes already occupied, Chapter~55's stimulated-emission secret behind lasers.

\step{Explaining the Opening Phenomenon}

Return to your palm. Overlapping skin and table clouds confront already occupied states. Exclusion forces some electrons into higher empty states, increasing energy under compression and producing pressure through its volume derivative. The wall combines familiar electromagnetic energy changes and a quantum constraint, no new force. Water molecules already sit near equilibrium under hydrogen bonds and van der Waals attraction; compressing further overlaps clouds and costs degeneracy energy. Their \(2.2\times10^{9}\)~Pa bulk modulus dwarfs hand pressure. Air molecules sit roughly ten molecular diameters apart without overlapping clouds; compression mostly removes empty gaps at thermal-energy cost. Nearly empty atoms give different mechanics at different densities.

Neon and sodium also settle. Neon, \(Z=10\), fills 2p; another electron must open shell three, feeling effective charge near \(1\), unfavorable addition. Removing a closed-shell electron costs \(21.6\)~eV. Blocked both ways, neon rests. Sodium, \(Z=11\), sends one extra electron to shell three, screened by ten inner electrons to roughly \(1.8\) effective charge and removable for \(5.14\)~eV. Losing it exposes a neon-like core, making electron disposal highly favorable. Water's oxygen accepts sodium's bargain, and released heat ignites hydrogen. Every periodic row closes with a full-shell idler and begins with a lone-electron hothead: periodicity marks energy filling, not a calendar.

Salt's yellow flame promotes the lone 3s electron to 3p. Its return emits about \(2.1\)~eV at \(589\)~nm. Exclusion's third-floor exile announces itself cheaply, completing the optional experiment's account and anticipating Chapter~55's spectroscopy.

\begin{figure}[htbp]
\centering
\begin{tikzpicture}[>=Stealth,scale=1.0]
  % Nucleus
  \draw[thick,fill=red!30] (0,0) circle (0.28);
  \node at (0,0) {\scriptsize \(+11\)};
  % Inner electron cloud
  \draw[thick,dashed,fill=blue!8] (0,0) circle (0.95);
  \node at (0,-1.25) {\scriptsize 10 inner electrons: \(-10\)};
  % Outer electron
  \draw[fill=black] (3.2,0) circle (0.06);
  \node[above] at (3.2,0.1) {\scriptsize 3s electron};
  % Effective-charge arrow
  \draw[->,very thick,orange!80!black] (0.35,0) -- (2.9,0);
  \node[above,orange!80!black] at (1.6,0.65) {\scriptsize \shortstack{Effective attraction\\about \(+1.8\)}};
\end{tikzpicture}
\caption{Screening: sodium's \(+11\) nucleus is screened by roughly \(-10\) inner-cloud charge, leaving its 3s electron an effective \(+1.8\), among the periodic table's easiest electrons to remove.}
\end{figure}

\step{Thought Experiments and Challenges}

\begin{thought}[If Electrons Were Bosons]
Change one cosmic rule: unlimited electrons per state. Within minutes all atomic electrons collapse into 1s, radii shrink as \(1/Z\), uranium ninety-two times smaller than hydrogen. Shells, closed shells, and periodicity vanish. Without exclusion walls, bonding is rewritten and solids lose hardness; interpenetrating clouds may erase even the distinction between hand and table. Push to stars: after massive stars exhaust fuel, gravity crushes atoms, packing electrons near \(10^{36}\)~m\(^{-3}\) and Fermi energies around \(10^5\)~eV, a hundred thousand chemistry scales. Exclusion's pressure supports Earth-sized, solar-mass white dwarfs. Above roughly \(1.4\) solar masses even this fails; electrons enter protons to become neutrons, whose degeneracy pressure supports neutron stars. Your table-pressing hand and white-dwarf stability obey one quantum account.
\end{thought}

\begin{puzzles}
\item Use exchange sign reversal, not formulas, to show electrons cannot share all four quantum numbers yet can share position. Why no contradiction? \emph{Hint}: exclusion constrains states, not positions; opposite spins can occupy spatially coincident probability distributions. Separating state from position is this chapter's key conceptual sorting.
\item Why does potassium's \(Z=19\) nineteenth electron occupy 4s rather than 3d? What effective charge would a forced 3d electron see? \emph{Hint}: argon's eighteen-electron core almost completely screens 3d, leaving \(Z_{\mathrm{eff}}\) near \(1\). Penetrating 4s feels appreciably more than \(1\), lowering its energy. Screening, not arbitrary rules, causes queue-jumping.
\item Count helium-3 and helium-4 fermions to explain similar chemistry but different low-temperature behavior. What fixes composite statistics? Is carbon-12, six protons, six neutrons, six electrons, fermionic or bosonic? \emph{Hint}: even constituent-fermion counts give bosons, odd counts fermions. Exchanging composites exchanges every constituent, so only parity sets the sign.
\item Someone says Pauli creates a strong electron repulsion force, making tables hard. Identify the switched concept and restate causality through energy versus volume. \emph{Hint}: exclusion is a state constraint without a force formula. Compression raises occupied energies, so \(E(V)\) rises steeply as volume shrinks, giving pressure \(-\mathrm dE/\mathrm dV\). Force is an energy slope, not a new force species, like the centripetal-force lesson.
\end{puzzles}
